% Reciprocal proofs for the separate Remark3.11(iii) integration fragment.
\section{Proofs for the ordered fixed-arity marker extension}

\begin{annexitem}{annex:p12-klop1987-om-bound}{thm:p12-klop1987-om-bound}{The actual reachable-signature bound}
\begin{proof}
Erase stars when reading numerical labels. PUT and PUSH change only
star bits; SELECT deletes a root and some sibling subtrees; LEX keeps
the same root label and copies only subtrees already present. COPY
copies its input and introduces only its new root \(m<n\). Thus no
operation introduces a number larger than the largest number in its
input. Context closure preserves the same assertion because unchanged
context nodes keep their numbers. Induction on the number of steps
proves the bound for every finite prefix of any run. In an infinite
run every node at every stage belongs to some finite prefix, so the
same bound holds throughout.

There is only one symbol of each label in the source signature; hence
\(\Sigma_M\subseteq\{0,\ldots,M\}\) has at most \(M+1\)
members. Each has one fixed finite arity. These facts give exactly a
finite ranked alphabet. They give no size or depth bound: COPY and
LEX can repeatedly increase both. The proof does not infer finiteness
of the set of reachable terms from finiteness of the alphabet.
\end{proof}
\end{annexitem}

\begin{annexitem}{annex:p12-klop1987-om-sets}{thm:p12-klop1987-om-sets}{Nonempty lower sets and finite witness domination}
\begin{proof}
The signature \(\Sigma_M\cup\{\bot\}\) is finite and ranked;
its precedence is strict and total, and its status is the fixed
left-to-right permutation. These are the precise hypotheses of the
LPO structural, finite-tree and well-foundedness results. Their full
proofs give a transitive irreflexive order, proper-subterm and context
closure, totality on distinct ground terms, and well-foundedness.
In particular, \(n(\vec b)>\bot\) for every actual source symbol
\(n\): use the strictly greater head and the empty list of target
arguments in the LPO root clause. Induction on terms gives that
\(\bot\) is below every other ground term in this semantic signature.

Both \(D(v)\) and \(D^-(v)\) are lower sets, by transitivity.
The recursive definition of \(J\) uses only the finitely many proper
children of its argument, so it is a well-defined structural recursion.
Prove its stated properties by induction on tree size. Each child
interpretation contains \(\bot\) by induction; hence the product of
child interpretations contains the all-bottom tuple. For arity zero
it instead contains its unique empty tuple. The resulting application
\(n(\vec b)\) is greater than \(\bot\), so both its weak and
strict lower sets contain \(\bot\). Their unions therefore are
nonempty and contain \(\bot\), including at a marked nullary node.
A union of lower sets is a lower set, proving that part as well.

If \(B_i\subseteq A_i\), every tuple used in the constructor
formed from the \(B_i\)'s is already a tuple used for the \(A_i\)'s.
This proves monotonicity for the plain and marked constructors alike.
Induct on the number of ancestors of a context hole, applying this
monotonicity at each ancestor. Stars on those ancestors do not change
the argument. Thus interpretation inclusion is preserved in all
finite contexts, not only the source's plain-root contexts.

Ground totality means that any two elements are equal or strictly
comparable. Repeatedly taking the larger of two proves that a finite
nonempty family has a greatest member which is one of its members.
If all members lie in \(A_i\), this maximum also lies in \(A_i\).
For finitely many witness tuples, take such a maximum in each
coordinate. Their common dominating tuple still belongs to
\(\prod_i A_i\). A coordinate without a stipulated witness may
be filled with \(\bot\). This constructs the finite maxima used
below; it does not assume completeness of LPO or a maximum for an
infinite argument set.

Finally if \(a_i\geq b_i\) in every coordinate, context closure
and transitivity imply \(n(\vec a)\geq n(\vec b)\): replace
strictly different coordinates successively, keeping equal ones.
If some coordinate is strictly different the result is strict.
This coordinate monotonicity is valid at every fixed finite arity.

To verify the stated failure case without bottom, use only a nullary
\(c\) and a binary \(n>c\) in the ground semantic alphabet.
Every ground term other than \(c\) has head \(n\), so is greater
than \(c\) by the strict head clause. Thus \(D(c)=\{c\}\)
and \(D^-(c)=\varnothing\). A marked \(c^*\) has empty
interpretation. Consequently the product of child interpretations
for \(n^*(c^*,c)\) is empty and its interpretation is empty.
SELECT of the second child has the nonempty interpretation
\(D(c)\). This is a literal failure of inclusion, not just an
unproved nonemptiness assumption. The actual interpretation above
uses the fresh lower \(\bot\), so no such child interpretation
is empty.
\end{proof}
\end{annexitem}

\begin{annexitem}{annex:p12-klop1987-om-operations}{thm:p12-klop1987-om-operations}{Every root case, including independent copies}
\begin{proof}
Write \(A_i=J(u_i)\) and
\(X=J(n^*(u_1,\ldots,u_k))\). It suffices to prove inclusion
for a root operation; the previous lemma transports it through any
context. All argument sets contain \(\bot\).

For PUT, each \(D^-(n(\vec a))\) is contained in
\(D(n(\vec a))\), so the union defining the marked result is
contained in that defining the plain input.

For SELECT of \(u_i\), take any \(b\in A_i\) and form a
tuple with coordinate \(i\) equal to \(b\) and all other
coordinates equal to \(\bot\). This is a permissible tuple.
The proper-subterm property gives \(n(\vec a)>b\), so
\(b\in X\). Hence \(J(u_i)\subseteq X\). The use of
bottom here is essential: no discarded argument makes the product
of argument choices empty, even if that argument has a marked leaf.

For COPY let \(m<n\), \(h=a(m)\). An arbitrary generator of
the plain right-hand interpretation is
\(m(b_1,\ldots,b_h)\), with each \(b_j\in X\).
When \(h>0\), choose for each \(j\) a witness tuple
\(\vec a^{\,j}\in\prod_i A_i\) satisfying
\(n(\vec a^{\,j})>b_j\). Such a witness exists by the
definition of \(X\). Take coordinatewise maxima of these finitely
many witness tuples to obtain \(\vec a\in\prod_i A_i\).
Coordinate monotonicity gives
\(n(\vec a)\geq n(\vec a^{\,j})>b_j\) for every \(j\).
The strict head comparison \(n>m\), together with this domination
of every target child, is exactly the second LPO clause; it yields
\(n(\vec a)>m(\vec b)\). Every member of the weak lower set
of this right-hand generator is therefore strictly below
\(n(\vec a)\), and belongs to \(X\). Taking their union
proves inclusion. If \(h=0\), choose any input tuple, for example
the all-bottom one. The strict head clause immediately gives
\(n(\vec a)>m()\), with no child obligations. If \(k=0\),
the input witness tuples throughout this calculation are the empty
tuple and their coordinatewise maximum is likewise empty. Thus
nullary inputs and zero-copy outputs are both covered.

For PUSH write the chosen plain-root child as
\(u_i=p(v_1,\ldots,v_h)\). A right-hand generator is
\(n(b_1,\ldots,b_k)\), with the unchanged \(b_j\in A_j\)
and \(b_i\in J(p^*(\vec v))\). By the latter membership
there is \(\vec c\in\prod_r J(v_r)\) with
\(p(\vec c)>b_i\). Put \(a_i=p(\vec c)\in J(u_i)\)
and \(a_j=b_j\) for \(j\ne i\). Context closure gives
\(n(\vec a)>n(\vec b)\). Every term weakly below the
right-hand generator is therefore strictly below this input witness
and lies in \(X\). Union over generators proves inclusion.
This argument explicitly uses that the pushed child has a plain
root, so that \(p(\vec c)\) belongs to its weak interpretation.

For LEX let \(k\geq1\) and write its plain-root first child
as \(u_1=p(\vec v)\). A right-hand generator is
\(n(b_1,b_2,\ldots,b_k)\), where
\(b_1\in J(u_1^*)\), while \(b_j\in X\) for \(j\geq2\).
Choose a witness \(a_1^0=p(\vec c)\in A_1\) with
\(a_1^0>b_1\). For each \(j\geq2\), choose a possibly
different tuple \(\vec a^{\,j}\in\prod_i A_i\) with
\(n(\vec a^{\,j})>b_j\). Take the first-coordinate maximum
of \(a_1^0\) and all \(a_1^{\,j}\), and in every other
coordinate take the maximum of the \(a_i^{\,j}\)'s; when no
such latter witness occurs, use \(\bot\). The resulting
\(\vec a\) is in \(\prod_i A_i\). It satisfies
\(a_1\geq a_1^0>b_1\), and, by coordinate monotonicity,
\(n(\vec a)>b_j\) for all \(j\geq2\). Also
\(n(\vec a)>a_1>b_1\), by the proper-subterm property.
Thus the source dominates every target argument, and its first
argument is strictly larger than the target's first argument.
The equal-head left-to-right LPO clause gives
\(n(\vec a)>n(b_1,\ldots,b_k)\). The right-hand generator's
entire weak lower set lies in \(X\), as required. At \(k=1\)
there are no copy witnesses and the same calculation uses just
\(a_1=a_1^0>b_1\). This proves the LEX case at every arity,
without requiring different copies to select the same values.

All five root cases are proved. The all-context monotonicity lemma
then gives inclusion for every auxiliary step, and transitivity of
set inclusion gives it for every finite auxiliary reduction.
\end{proof}
\end{annexitem}

\begin{annexitem}{annex:p12-klop1987-om-first-mark}{thm:p12-klop1987-om-first-mark}{Strictness is supplied by the first PUT, not assumed for all steps}
\begin{proof}
Induct on a plain term \(t=n(t_1,\ldots,t_k)\). By induction
\(J(t_i)=D(t_i)\). Every tuple \(\vec b\) in these child
interpretations satisfies \(t_i\geq b_i\), hence
\(t\geq n(\vec b)\) by coordinate monotonicity. Thus the
union defining \(J(t)\) is contained in \(D(t)\). The
particular tuple \((t_1,\ldots,t_k)\) belongs to that product
and contributes all of \(D(t)\); therefore equality holds.
For a root-marked \(t^*\), the same reasoning shows that every
strict lower set used in its interpretation is contained in
\(D^-(t)\), and the original tuple contributes all of
\(D^-(t)\). The argument includes nullary terms.

Prove strict lifting up a plain context by induction on its depth.
At the hole, \(J(v^*)=D^-(v)\). Suppose the interpretation of
the changed child is contained in \(D^-(v)\), while its plain
siblings \(t_j\) have interpretations \(D(t_j)\). In any
generator of the enclosing plain constructor, the changed value
\(b_i\) satisfies \(v>b_i\), and the unchanged chosen values
satisfy \(t_j\geq b_j\). Replacing the changed coordinate
strictly and every other differing coordinate weakly gives
\(n(t_1,\ldots,v,\ldots,t_k)>n(\vec b)\).
The entire weak lower set of the generator is strictly below the
original parent. Taking the union proves the same strict-lifting
property at that parent; iteration proves it at the whole context.

A positive reduction from an entirely plain tree begins with PUT:
every other root rule requires an existing mark, and no such mark
exists anywhere in the initial tree. After this first step its
interpretation is contained in \(D^-(t)\), by strict lifting.
Every later step preserves inclusion by the operations lemma, so
the final interpretation remains contained in \(D^-(t)\).
If the final tree \(s\) is plain, then \(s\in J(s)=D(s)\).
Consequently \(s\in D^-(t)\), or \(t>s\), proving the endpoint
claim. This also explains why merely assigning nonincreasing
interpretations to operations would have been insufficient without
the strict first-mark calculation.
\end{proof}
\end{annexitem}

\begin{annexitem}{annex:p12-klop1987-om-well-founded}{thm:p12-klop1987-om-well-founded}{A global theorem, using an actual invariant rather than a restricted input class}
\begin{proof}
The positive relation on plain endpoints is transitive by concatenating
finite reductions whose common endpoint is plain. A positive cycle
would imply \(t>t\) by the first-mark lemma in the finite signature
for that cycle. LPO irreflexivity excludes it.

Suppose there were an infinite sequence of plain trees
\(t_0\rightsquigarrow_{\rm lex}^{+}t_1
\rightsquigarrow_{\rm lex}^{+}t_2\rightsquigarrow_{\rm lex}^{+}\cdots\).
Choose a finite witnessing reduction for each displayed edge.
All witnesses concatenate to a run starting at the finite tree
\(t_0\). Let \(M\) be its largest numerical label. The
reachable-signature lemma shows that every term in every witness
lies over the same finite \(\Sigma_M\). Use the single semantic
alphabet \(\Sigma_M\cup\{\bot\}\), precedence and LPO for
all edges. The first-mark lemma then gives
\(t_0>t_1>t_2>\cdots\), contradicting the proved finite-ranked
LPO well-foundedness. It is not legitimate to choose an unrelated
finite order for each edge; the invariant supplies one for the
entire hypothetical chain.

This proves well-foundedness for arbitrary ambient
\(\Sigma\subseteq\mathbb N\) and arbitrary fixed finite
arities. The initial \(M\) depends on the individual chain's
first term; it is not a common bound imposed on all source objects.
No argument bounds term depth, and the conclusion also holds for
the source subrelation of the all-context auxiliary game.

For completeness the marked cycle displayed in the body is well
typed: LEX at the marked binary root requires only the plain root
of \(c\); PUT next marks the new plain binary root; SELECT next
chooses its second child, which is literally the original
\(n^*(c,u)\). Its marked endpoints fall outside the proved
restriction. Thus no stronger auxiliary-SN assertion follows.
\end{proof}
\end{annexitem}

\begin{annexitem}{annex:p12-klop1987-om-applications}{thm:p12-klop1987-om-applications}{Actual positive paths and ordinary open-term termination}
\begin{proof}
Use the labels and arities stated for \(A,S,0\). For the first
rule the positive path is
\[
 A(0,x)\longrightarrow A^*(0,x)
 \longrightarrow S(A^*(0,x))\longrightarrow S(x).
\]
These are PUT, COPY to the smaller unary \(S\), and SELECT of
the second child of the marked \(A\), in a plain \(S\)-context.
For the second rule a full path is
\begin{align*}
 A(S(x),0)&\longrightarrow A^*(S(x),0)\\
 &\longrightarrow A(S^*(x),A^*(S(x),0))\\
 &\longrightarrow A(x,A^*(S(x),0))\\
 &\longrightarrow A(x,S(A^*(S(x),0)))\\
 &\longrightarrow A(x,S(0)).
\end{align*}
The operations are PUT, LEX, SELECT of the unary marked \(S\),
COPY of the remaining marked \(A\) to unary \(S\), and SELECT
of its second child. For the third rule the report's exact path is
\begin{align*}
 A(S(x),S(y))&\longrightarrow A^*(S(x),S(y))\\
 &\longrightarrow A(S^*(x),A^*(S(x),S(y)))\\
 &\longrightarrow A(x,A^*(S(x),S(y)))\\
 &\longrightarrow A(x,A(S(x),S^*(y)))\\
 &\longrightarrow A(x,A(S(x),y)).
\end{align*}
The last two operations PUSH into the second argument of that
marked \(A\), then SELECT the only child of the marked \(S\).
Each pushed or lexically marked child has the required plain root.
The symbols \(x,y\) in these displays are metavariables for
arbitrary plain trees; no numerical-normality premise is used.

For associativity write \(d\) for its one binary operator. The
path, at its one fixed arity and without sibling permutations, is
\begin{align*}
 d(d(x,y),z)&\longrightarrow d^*(d(x,y),z)\\
 &\longrightarrow d(d^*(x,y),d^*(d(x,y),z))\\
 &\longrightarrow d(x,d^*(d(x,y),z))\\
 &\longrightarrow d(x,d(d^*(x,y),z))\\
 &\longrightarrow d(x,d(y,z)).
\end{align*}
These steps are PUT, LEX, SELECT of the first marked inner
\(d\)'s first child, PUSH into the other marked \(d\)'s first
child, and SELECT of that child's second argument. No commutative
tree identification is involved.

To pass from these tree paths to all ordinary finite terms, encode
the original function symbols by the specified numerical labels.
For the Ackermann system encode every ordinary object variable by
one additional nullary symbol \(c\) with label3. For associativity
give \(d\) label1 and every variable the same fresh nullary label0.
This homomorphic coding need not distinguish different variables.
It sends every ordinary finite term to a plain tree in one fixed
ranked numerical signature. Any instantiated rule becomes exactly
one of the positive paths just proved, with its metavariables
replaced by the coded instance terms. Use of a rule in an ordinary
context becomes the same path in its coded plain context. Consequently
every ordinary step gives a positive plain-to-plain marker reduction.
If two coded endpoints happened to coincide, that path would itself
be a positive plain cycle, already excluded; injectivity is not a
missing premise of this transport.

An infinite ordinary reduction would therefore give an infinite
sequence of positive plain-endpoint marker reductions, contradicting
the global theorem. This proves SN on open and ground finite terms
for all three source Ackermann rules and for the source's stated
associativity orientation. It proves neither the opposite orientation
in this same left-to-right status nor unrestricted commutativity.
\end{proof}
\end{annexitem}
